\scanpage{59}
\mspage{8}
% archive sheet number: 59

moyennant les axiomes qu'il faut sur $W$. J'ai envie de regarder d'abord le choix $S'=\underline{\mathrm{Ch}}(S;s_0,-)$, sans passage à une limite inductive qui m'obligerait à mettre un axiome de limites sur $W$. Les \deleted{\unsure{théorie}} propriétés homotopiques délicates et essentielles de $\underline{\mathrm{Ch}}(S)$ sont établies dans $\overline{\underline{\mathrm{VII}}}$, p. 157-169, savoir que dans le diagramme

\diagram{\begin{tikzcd}[ampersand replacement=\&]
S\arrow[r,"i_S",dash] \& \underline{\mathrm{Ch}}(S)\arrow[r,"p_S"] \& S\times S
\end{tikzcd}}
, $i_S\in W_\infty$ et $p_S$ est une \deleted{\unsure{$W$}}$\unsure{\mathrm{Hot}|W_\infty}$-fibration (en fait, \uline{propre} et \uline{lisse}\inserted{ pour $W_\infty$}). Le fait que $i_S\in W$ implique que les \unsure{foncteurs} source et but

% slanted
\margin{Dans l'énoncé des ``propriétés ess.'' de $\underline{\mathrm{Ch}}(S)$, on peut travailler avec $W_\omega$, $i_S\in$ la plus petite partie de $\mathrm{Fl}(\mathrm{Cat})$ satisfaisant les axiomes $W(1)$, $W(2\,\mathrm{bis})$, $W(3)$ (cf \unsure{p.ge|page} 14), donc $W_\omega\subset W\cap W_\infty$}

\diagram{\begin{tikzcd}[ampersand replacement=\&,column sep=large]
\underline{\mathrm{Ch}}(S)\arrow[r,shift left=.6ex,"\sigma",dash]\arrow[r,shift right=.6ex,"\beta"',dash] \& S
\end{tikzcd}}
\[
\sigma=\mathrm{pr}_1\circ p_S,\quad \beta=\mathrm{pr}_2\circ p_S
\]

sont dans $W_\infty$ (puisque $\sigma i_S=\beta i_S=\mathrm{id}_S$). Comme \deleted{c'est une} sont des $\mathrm{Hot}$-fibrations, il en résulte que les fibres $\underline{\mathrm{Ch}}(S;s,-)$ ou $\underline{\mathrm{Ch}}(S;-,s)$ sont $\unsure{\mathrm{Hot}}$-asphériques. Mais même en admettant que $\sigma,\beta$ sont dans $W$ et sont des \deleted{\unsure{$\mathrm{Hot}$}}$W$-fibrations, je ne peux en conclure \inserted{tout de suite} que les fibres sont $W$-asphériques -- il y aurait cercle vicieux ! Il me faut donc une démonstration directe.

Considérons le foncteur
\diagram{\begin{tikzcd}[ampersand replacement=\&,row sep=small]
\underline{\mathrm{Ch}}(S;s,-)\arrow[d,"p"']\\
\underline{\mathrm{Ty}}
\end{tikzcd}}

, c'est un foncteur $\mathrm{Cat}$-fibrant, donc $W$-lisse (je \deleted{devrai} \unsure{développer} la notion de $W$-lissité \unsure{bien} sûr, \deleted{\unsure{je dois}} \unsure{j'utiliserai|j'admettrai} pour le moment les propriétés formelles familières, et les justifierai par la suite). A cause de la $W$-lissité, pour prouver $p\in W$, il suffit de

\margin{\deleted{\unsure{$f$}}$f:X\unsure{\text{---}|\longrightarrow}Y$ est $W$-lisse si les $X_y\longrightarrow{}_{y\backslash}X$ sont $W$-asphériques}
